\section{Introduction}
\label{sec:intro}

As the internet becomes saturated with content generated by large language models (LLMs), future models will inevitably be trained on AI-generated data. There is growing interest in understanding the long-term sustainability of this cycle. Prior research indicates that when a model is recursively trained on its own outputs, it suffers from ``model collapse''---a phenomenon where the generated distribution loses diversity, rare events disappear, and the model converges to a homogenized state \cite{shumailov2024ai}.

Given the ability of LLMs to generate vast amounts of data, identifying the conditions under which an ecosystem of models can survive is critical. Biological ecosystems rely on a minimum viable population to avoid extinction, prompting a parallel question in artificial intelligence: is there a minimum viable population for an LLM information ecosystem such that it does not collapse?

While previous studies primarily focus on single-model recursive loops \cite{shumailov2024ai}, there is a significant gap in defining the quantitative threshold or the exact nature of ``individuality'' required to sustain a multi-agent population. However, it also becomes increasingly challenging to disentangle the effect of sheer numbers from the effect of structural diversity. A population of identical models might simply reinforce shared biases faster, acting as an echo chamber rather than a stabilizing force.

We hypothesize that an ecosystem's viability depends more on epistemic diversity than on the absolute number of agents. To investigate this, we empirically test the ``Minimum Viable Population'' (MVP) hypothesis by simulating ecosystems of varying sizes and diversity levels. We simulate a recursive training ecosystem using small LLMs, comparing a single-model baseline against both homogeneous and heterogeneous multi-model populations.

Our key findings demonstrate that population size alone is insufficient to prevent collapse. In fact, our homogeneous population of three models exhibits a more severe drop in unique sentence ratio compared to the single-model baseline (47.3\% vs. 73.5\%).

We provide the following contributions:
\begin{itemize}[leftmargin=*,itemsep=0pt,topsep=0pt]
    \item We propose a multi-agent recursive training framework to empirically test the minimum viable population hypothesis for LLMs.
    \item We demonstrate that increasing the population size of identical models sharing a data pool accelerates model collapse.
    \item We show that introducing epistemic diversity via different sampling temperatures mitigates collapse, improving ecosystem diversity over homogeneous populations.
\end{itemize}
